\section{Diseño conceptual mecatrónico}

\subsection{Descripción de la solución propuesta}
El diseño del sistema CIP se ha concebido como una solución compacta e integrada, orientada a optimizar los procesos de limpieza y desinfección de tuberías mediante una secuencia automatizada de etapas. La arquitectura funcional del equipo se articula en torno a tres subsistemas de almacenamiento: un tanque para la recogida de agua residual —cuyo llenado se realiza por gravedad al término de cada ciclo—, un tanque destinado al agente de limpieza y desinfección, y una cabina de operaciones que constituye el núcleo reactivo del sistema, donde se acondicionan los fluidos antes de su descarga a la línea objetivo.

El proceso productivo se estructura en tres ciclos diferenciados, ejecutados de manera secuencial dentro de la cabina de operaciones:

Ciclo I (preenjuague térmico): se procede al llenado de la cabina con agua hasta alcanzar un nivel preestablecido, monitorizado por sensor. Acto seguido, se activa la resistencia de inmersión hasta lograr la temperatura de consigna, tras lo cual el volumen acondicionado es evacuado hacia la tubería a limpiar.

Ciclo II (lavado químico y mecánico): nuevamente se carga la cabina con agua hasta el nivel de referencia. En esta etapa se incorpora una dosificación controlada del agente jabonoso, al tiempo que se activa el calentamiento y se acciona el agitador (motoreductor) para asegurar la homogeneización y la eficacia térmica de la solución. Finalizada la preparación, la mezcla es impulsada hacia la tubería.

Ciclo III (enjuague final): se realiza un último llenado con agua hasta el nivel establecido, para ser descargado directamente sin tratamiento térmico ni aditivos, garantizando el arrastre de los residuos remanentes y concluyendo la secuencia de limpieza.

El abastecimiento de agua proviene de la red general, cuyo ingreso a la cabina es regulado mediante una electroválvula de paso, gobernada por la lógica de control, que habilita o interrumpe el flujo según la demanda de cada ciclo.

En cuanto al manejo de fluidos, el sistema dispone de dos bombas con cometidos claramente diferenciados: una destinada a la extracción de los insumos procesados desde la cabina hacia la tubería, y otra encargada de la dosificación precisa del jabón desde su tanque de almacenamiento. Con el fin de garantizar la estanqueidad y evitar el reflujo o la contaminación cruzada entre fluidos, a la salida de cada bomba se ha dispuesto una configuración de doble bloqueo mediante dos electroválvulas en serie, que aseguran el aislamiento hidráulico cuando la bomba no se encuentra en operación.

El acondicionamiento térmico del agua se efectúa exclusivamente a través de una resistencia de inmersión, cuya potencia de $2000\,\mathrm{W}$ permite alcanzar los rangos de temperatura requeridos en los tiempos establecidos.

Finalmente, la instrumentación de seguridad incorpora sensores de nivel tanto en el tanque de agua residual como en la cabina de operaciones, cuyas señales permiten prevenir condiciones de rebose y garantizar la integridad operativa del conjunto, deteniendo el llenado de forma automática al alcanzar los límites máximos permitidos.
\insertarfigura[0.5\textwidth]{GENERAL CIP.jpeg}{CAD general del sistema CIP.}{fig:general}

\subsection{Diagrama general de bloques}
\insertarfigura{diagrama_general.pdf}{Diagrama general del sistema.}{fig:general-bloques}

Esta configuración responde directamente a las especificaciones 
técnicas del proyecto. La integración del agitador AG-101 resuelve 
la exigencia de homogeneidad y valida la hipótesis de temperatura 
uniforme en el tanque. El lazo de control entre el sensor TT-101 
y la válvula SV-101 asegura la precisión térmica. Por su parte, 
la bomba P-102 garantiza el dinamismo hidráulico para la limpieza 
de la tubería, mientras que la disposición final en TK-103 o 
desagüe cumple con el manejo seguro de efluentes en paso directo.

\subsection{Diagrama general de bloques}
\insertarfigura{diagrama_general.pdf}{Diagrama general del sistema.}{fig:diagrama}
\subsection{Arquitectura funcional}

La arquitectura funcional del sistema CIP organiza la operación como una secuencia controlada de etapas, desde la preparación de la solución hasta la descarga final. En cada etapa, la lógica de control coordina el ingreso de agua, la dosificación de jabón, la agitación, el calentamiento, la impulsión del fluido y la supervisión de los niveles para garantizar una limpieza segura y repetible de la tubería
esto se desea mostrar en una CadeSimu el cual permite ver la integracion de un Grafcet preliminar no sera el usado de manera definiticva pero permite ver el funcionamiento del sistema.

\insertarfigura[0.6\textwidth]{Cade_Pi.png}{Arquitectura funcional del sistema.}{fig:arquitectura}

\begin{figure}[H]
\centering
\begin{minipage}[t]{0.47\textwidth}
\raggedright\textbf{Etapa 1. El proceso inicia al darle la señal de arranque el cual permite el accionamiento de la electrovalvula 1 la cual permite el llenado de la cabina con agua}\\[2mm]
\IfFileExists{figuras/etapa-01.png}{\includegraphics[width=\textwidth]{etapa-01.png}}{\fbox{\parbox[c][4cm][c]{1.0\textwidth}{\centering Inserte la imagen de la etapa 1}}}
\captionof{figure}{Etapa 1: Llenado de la cabina de operaciones.}\label{fig:etapa-01}
\end{minipage}\hfill
\begin{minipage}[t]{0.47\textwidth}
\raggedright\textbf{Etapa 2. Al llegar al nivel de referencia se acciona la resistencia de inmersión para calentamiento del agua hasta una temperatura determinada por el controlador PID}\\[2mm]
\IfFileExists{figuras/etapa-02.png}{\includegraphics[width=\textwidth]{etapa-02.png}}{\fbox{\parbox[c][4cm][c]{0.9\textwidth}{\centering Inserte la imagen de la etapa 2}}}
\captionof{figure}{Etapa 2: Calentamiento del agua mediante resistencia de inmersión.}\label{fig:etapa-02}
\end{minipage}
\end{figure}

\begin{figure}[H]
\centering
\begin{minipage}[t]{0.47\textwidth}
\raggedright\textbf{Etapa 3. Al cumplirse la condición de calentamiento se acciona la bomba de drenado de la cabina de operacion llevado el agua de prejuagado a la tuberia a limpiar (esto solo puede suceder si el tanque de agua resifual se encuentra en el nivel optimo para no desbordar y generar derramamiento)}\\[2mm]
\IfFileExists{figuras/etapa-03.png}{\includegraphics[width=\textwidth]{etapa-03.png}}{\fbox{\parbox[c][4cm][c]{0.9\textwidth}{\centering Inserte la imagen de la etapa 3}}}
\captionof{figure}{Etapa 3: Bomba de drenado de la cabina de operaciones.}\label{fig:etapa-03}
\end{minipage}\hfill
\begin{minipage}[t]{0.47\textwidth}
\raggedright\textbf{Etapa 4. se realiza el rellenado de la cabina de operaciones esto se realizara cuando ya el nivel del tanque se detecta como bajo}\\[2mm]
\IfFileExists{figuras/etapa-04.png}{\includegraphics[width=\textwidth]{etapa-04.png}}{\fbox{\parbox[c][4cm][c]{0.9\textwidth}{\centering Inserte la imagen de la etapa 4}}}
\captionof{figure}{Etapa 4: Rellenado de la cabina de operaciones.}\label{fig:etapa-04}
\end{minipage}
\end{figure}

\begin{figure}[H]
\centering
\begin{minipage}[t]{0.47\textwidth}
\raggedright\textbf{Etapa 5. Dosificacion del jabon, al estar lleno el tanque se accionara la bomba de dosificación del la sustancia de limpieza}\\[2mm]
\IfFileExists{figuras/etapa-05.png}{\includegraphics[width=\textwidth]{etapa-05.png}}{\fbox{\parbox[c][4cm][c]{0.9\textwidth}{\centering Inserte la imagen de la etapa 5}}}
\captionof{figure}{Etapa 5: Dosificación del jabón.}\label{fig:etapa-05}
\end{minipage}\hfill
\begin{minipage}[t]{0.47\textwidth}
\raggedright\textbf{Etapa 6. Se realiza la agitacion y el calentamiento de la mezcla de agua y sustancia de limpieza esto se hara luego de se tenga en cuentra que se estuvo accionado un tiempo determinado la bomba de dosificacion de sustencia de limpieza}\\[2mm]
\IfFileExists{figuras/etapa-06.png}{\includegraphics[width=\textwidth]{etapa-06.png}}{\fbox{\parbox[c][4cm][c]{0.9\textwidth}{\centering Inserte la imagen de la etapa 6}}}
\captionof{figure}{Etapa 6: Agitación y calentamiento de la mezcla de agua y sustancia de limpieza.}\label{fig:etapa-06}
\end{minipage}
\end{figure}
\begin{figure}[H]
\centering
\begin{minipage}[t]{0.47\textwidth}
\raggedright\textbf{Etapa 7. Se realiza el proceso de calentar y agitar durante un tiempo determinado y si el tanque de agua residual se encuentra en un nivel aceptable se realiza la operación de drenado de la cabina de operación de nuevo a la tubería a limpiar}\\[2mm]
\IfFileExists{figuras/etapa-07.png}{\includegraphics[width=\textwidth]{etapa-07.png}}{\fbox{\parbox[c][4cm][c]{0.9\textwidth}{\centering Inserte la imagen de la etapa 7}}}
\captionof{figure}{Etapa 7: circulación de la solución de limpieza en la tubería.}\label{fig:etapa-07}
\end{minipage}\hfill
\begin{minipage}[t]{0.47\textwidth}
\raggedright\textbf{Etapa 8. luego de ya drenar la mezcla si el nivel del tanque ya que se encuentra en un nivel aceptable se realiza la operacion de rellenado del tanque para el ultimo ciclo}\\[2mm]
\IfFileExists{figuras/etapa-08.png}{\includegraphics[width=\textwidth]{etapa-08.png}}{\fbox{\parbox[c][4cm][c]{0.9\textwidth}{\centering Inserte la imagen de la etapa 8}}}
\captionof{figure}{Etapa 8: rellenado del tanque para el ultimo ciclo.}\label{fig:etapa-08}
\end{minipage}
\end{figure}
\begin{figure}[H]
\centering
\begin{minipage}[t]{0.47\textwidth}
\raggedright\textbf{Etapa 9. Se realiza el proceso de enjuague con agua para eliminar residuos remanentes}\\[2mm]
\IfFileExists{figuras/etapa-09.png}{\includegraphics[width=\textwidth]{etapa-09.png}}{\fbox{\parbox[c][4cm][c]{0.9\textwidth}{\centering Inserte la imagen de la etapa 9}}}
\captionof{figure}{Etapa 9: enjuague con agua para eliminar residuos remanentes.}\label{fig:etapa-09}
\end{minipage}\hfill
\end{figure}
\subsection{P\&ID conceptual y recorrido del fluido}
\insertarfigura{pid_conceptual (2).pdf}{P\&ID conceptual del sistema CIP.}{fig:pid}
\begin{table}[H]
\centering
\scriptsize
\caption{Identificación de elementos del diagrama P\&ID (Equipos, Válvulas, Instrumentos y Líneas).}
\label{tab:pid_elementos}
\begin{tabularx}{\textwidth}{L{2.5cm}L{3.5cm}L{3.0cm}X}
\toprule
\textbf{Tag / Identificador} & \textbf{Nombre} & \textbf{Categoría / Tipo} & \textbf{Función} \\ \midrule
TK-101 & Tanque de jabón & Equipo (Tanque) & Almacena la materia prima (jabón) antes de la dosificación. \\
P-101 & Bomba de jabón & Equipo (Bomba) & Impulsa el jabón hacia el tanque de mezcla. \\
TK-102 & Tanque de mezcla & Equipo (Tanque) & Recipiente principal donde se realiza la mezcla y calentamiento. \\
AG-101 & Agitador & Equipo (Interno) & Agita e homogeniza la mezcla dentro de TK-102. \\
M-101 & Motorreductor & Equipo (Motor) & Acciona mecánicamente el eje del agitador AG-101. \\
E-101 & Resistencia eléctrica calefactora & Equipo (Calefactor) & Proporciona energía térmica para calentar el agua/mezcla. \\
P-102 & Bomba de limpieza & Equipo (Bomba) & Impulsa la mezcla de limpieza hacia el circuito. \\
TK-103 & Tanque de almacenamiento de agua & Equipo (Tanque) & Almacena el efluente/agua de retorno tras el proceso de limpieza. \\
SV-101 & Válvula solenoide de jabón & Válvula de control & Permite o interrumpe el paso de jabón desde TK-101 a P-101. \\
SV-102 & Válvula solenoide de agua & Válvula de control & Regula el ingreso de agua de red hacia el tanque de mezcla. \\
SV-103 & Válvula solenoide de descarga & Válvula de control & Controla la salida de mezcla de TK-102 hacia la bomba P-102. \\
LIT-101 & Transmisor indicador de nivel & Instrumento & Mide e indica localmente el nivel de líquido en TK-102. \\
TT-101 & Transmisor de temperatura & Instrumento & Mide la temperatura del fluido dentro de TK-102. \\
L-1 & Línea de jabón & Tubería de proceso & Transporta jabón impulsado hacia la mezcla. \\
L-2 & Línea de agua & Tubería de proceso & Alimenta agua de proceso a TK-102. \\
L-3 & Línea de descarga de mezcla & Tubería de proceso & Salida de fluido desde la base de TK-102. \\
L-4 & Línea de efluente de limpieza & Tubería de proceso & Conduce el fluido que retorna de la tubería de limpieza. \\
L-5 & Línea de desagüe de agua & Tubería de proceso & Drena o evacúa el agua almacenada en TK-103. \\
--- & Señales de control (PLC) & Línea de instrumentación & Transmite señales eléctricas/datos de entrada y salida entre el PLC y los equipos/instrumentos (SV-101, SV-102, SV-103, LIT-101, TT-101, E-101, M-101, P-102). \\ \bottomrule
\end{tabularx}
\end{table}
\insertarfigura{recorrido_fluidos.pdf}{Recorrido preliminar de fluidos.}{fig:fluidos}

\subsection{Distribución preliminar y diseño CAD}
\insertarfigura{cad_prototipo.png}{Vista del diseño CAD preliminar.}{fig:cad}
\campo{Indique dimensiones, masa estimada, ubicación de tanque, bomba, calentador, válvulas y sensores, conexiones externas y criterios de portabilidad.}

\subsection{Equipo o circuito de prueba}
\campo{Identifique el equipo que será sometido al ciclo y describa sus conexiones de entrada, retorno y drenaje.}
